\section{Discussion}

The result refines rather than overturns the slingshot picture. Nicholson and Forgan's
speed-up is real and large \cite{nicholson-forgan-2013}, but it was computed under perfect
coordination, and coordination is not free. Under finite light-speed a meaningful fraction
of the advantage, roughly a third for nearest-star slingshots and half for maximum-boost,
is spent on trips to stars that turn out to be already taken. The more aggressively a
policy reaches for distant, high-boost targets, the more of its speed it loses to
uncoordinated collisions.

\subsection{Limitations}
The measurement is a conservative upper bound in two senses, and both point to future work.
Probes here learn only at their decision stars, ignoring news their paths cross in flight,
so a two-endpoint knowledge model would recover some of the wasted trips; and probes here
share nothing directly, so adding probe-to-probe relay of the settled map would let a swarm
partially coordinate rather than merely collide. Both would lower the tax without erasing
it, since the fundamental limit, that no probe can know a distant star is taken until its
light arrives, is set by physics, not by protocol. Two further simplifications bound the
scope. The field is at a uniform one star per cubic parsec; a realistic, clumpier
distribution would lengthen the long-range hops that carry the penalty, so the reported
figures are conservative for the maximum-boost policy in particular. And the dynamics have
no loss term: settled sites never fail, so the field always closes. A settlement-death term
of the kind that produces the ``Aurora'' steady state, in which the settled fraction
equilibrates below one \cite{carroll-nellenback-2019}, is exactly what would be needed to
turn lag into a permanent unsettled remainder; delay alone does not.

\subsection{The signature of delayed agreement}
The behaviour is the expected signature of delayed multi-agent agreement. Formal consensus
degrades as the ratio of communication delay to decision cadence grows: the standard
protocol is stable only while the one-hop delay stays below a bound set by the coupling
\cite{olfati-saber-murray-2004}, a result later generalized to switching topologies and
link failures \cite{olfati-saber-fax-murray-2007}. Our probes never reach agreement about
the settled map at all; at a mean inter-star spacing of about a parsec the round-trip
light-time is several years, far longer than any decision cadence, so they behave like
agents forced to act on purely local, out-of-date views. In the limit of unbounded delay,
guaranteed agreement is not merely slow but impossible \cite{flp-1985}, and a system that
must nonetheless keep acting is forced to trade a single consistent map for availability
\cite{gilbert-lynch-2002}, which is precisely the choice (\ref{eq:gate}) encodes.

This is the same progression that drove the space teleoperation and networking communities
away from real-time control. Past about a second of round-trip delay a human operator
abandons closed-loop control for move-and-wait \cite{ferrell-1965}, and as delay grows the
handoff continues through supervisory control to full remote autonomy \cite{sheridan-1993}.
Past minutes to hours no continuous end-to-end path exists at all, and communication becomes
hop-by-hop store-and-forward, the delay-tolerant regime motivated by the interplanetary case
\cite{burleigh-2003} and standardized first as an architecture \cite{rfc-4838} and now as a
deployed protocol \cite{rfc-9171}. At interstellar hop lengths every one of those faster
regimes is left far behind; each probe is effectively an independent colony acting on
priors, and the coordination tax is what that independence costs when probes want the same
stars.

\subsection{A design lesson}
There is a practical corollary. For a real swarm on stale maps, a locally greedy policy is
not merely simpler than a boost-maximising one; it is more robust to the delay it cannot
avoid, because a lost race costs only the star next door. A designer who cannot buy
coordination should prefer locality, spending the slingshot advantage on speed rather than
on reach.
